\section{AI 代码审查的新世界}

\subsection{主要工具}

\begin{itemize}
    \item \textbf{Graphite}：本周嘉宾 Tomas Reimers 是其 CPO
    \item \textbf{Greptile}：基于代码库理解的 AI 审查
    \item \textbf{CodeRabbit}：自动化 PR 审查
    \item \textbf{Claude Code / Codex}：通用 AI 编码工具也可用于审查
\end{itemize}

\subsection{AI 代码审查带来的改变}

\begin{enumerate}
    \item \textbf{效率提升}：自动化处理常规检查，缩短审查周期
    \item \textbf{一致性}：AI 对所有代码应用相同的标准
    \item \textbf{知识共享}：AI 可以携带整个代码库的上下文
    \item \textbf{减轻认知负担}：让人类专注于复杂逻辑和架构决策
    \item \textbf{持续改进}：AI 系统随时间积累更多数据而改善
    \item \textbf{整体理解}：AI 可以从全局视角审视变更的影响
\end{enumerate}

\begin{knowledgebox}{AI 代码审查的分工模式}
最佳实践不是用 AI 完全替代人类审查，而是让 AI 处理``表层''问题（格式、命名、常见 bug 模式、安全检查），让人类专注于``深层''问题（架构合理性、业务逻辑正确性、长期维护性）。
\end{knowledgebox}

\subsection{当前的局限}

\begin{itemize}
    \item \textbf{配置和设置成本}：需要初始投入来训练系统
    \item \textbf{假阳性}：需要持续训练系统以减少误报
    \item \textbf{无法捕获代码库的惯用法和最佳实践}：仓库特有的 idiom 仍然超出 AI 能力
    \item \textbf{无法处理复杂的业务逻辑和架构决策}：这正是人类仍然不可或缺的地方
    \item \textbf{安全变更需格外谨慎}：AI 可能遗漏微妙的安全影响
    \item \textbf{经常遗漏边界情况}
\end{itemize}

\begin{warningbox}{AI 时代代码审查更重要了，而不是更不重要}
当 AI 编写了你的大部分代码时，代码审查变得\textbf{更加}重要。你拥有被合并和部署的代码——不能把锅甩给 AI。``The AI wrote it'' 不是一个可接受的借口。
\end{warningbox}

\subsection{AI 审查系统的落地架构}

如果把 AI review 看成一个工程系统，而不是一个 chatbot，它至少包含以下流水线：

\begin{table}[H]
\centering
\begin{tabular}{p{0.18\textwidth} p{0.25\textwidth} p{0.4\textwidth}}
\toprule
阶段 & 关键输入 & 主要任务 \\
\midrule
变更解析 & Git diff、文件类型、提交信息 & 判断应该关注哪些模块、哪些变更具有高风险 \\
仓库检索 & 相关文件、历史 PR、代码规范 & 补足 diff 之外的上下文，减少脱离语境的评论 \\
规则执行 & Lint、安全规则、团队约束 & 过滤显而易见的问题，降低人工负担 \\
语义评估 & 逻辑流、接口契约、异常路径 & 发现更深层的正确性或维护性风险 \\
评论生成 & 风险等级、证据、建议方案 & 输出可执行、可讨论、可归因的 review comments \\
\bottomrule
\end{tabular}
\caption{AI 代码审查的典型流水线}
\end{table}

\begin{warningbox}{最大的失败模式是上下文装配错误}
很多 AI review 产品的问题并不是模型不会写评论，而是没有把足够的代码库上下文、团队规则和历史惯例喂给模型。上下文一旦错位，就会产生大量貌似合理但实际上没用的评论。
\end{warningbox}

\subsection{人类 reviewer 在哪里最不可替代}

课程和行业实践都指向同一个结论：AI 越擅长检查局部模式，人类 reviewer 越应该把时间投入到高阶判断。

\begin{itemize}
    \item \textbf{意图校验}：这次改动真的解决了正确的问题吗？
    \item \textbf{架构一致性}：它是否让系统朝更混乱的方向演化？
    \item \textbf{组织记忆}：是否违背了团队过往踩过坑后形成的隐性规则？
    \item \textbf{风险取舍}：某些缺陷是否可以接受，某些技术债是否应该延期？
\end{itemize}

\begin{importantbox}{未来的 reviewer 更像 ``最后签字的系统设计师''}
在 AI 生成代码的比例越来越高时，reviewer 不只是看格式或命名的人，而是对变更是否应该进入主干负责的人。这个角色要求更强的系统理解，而不是更快地扫 diff。
\end{importantbox}

\subsection{团队应该怎样衡量 review 系统是否真的变好}

如果没有度量，团队往往会把``评论数量变多''误判为``审查质量变高''。更合理的指标应同时覆盖速度、质量和信任。

\begin{itemize}
    \item \textbf{首轮反馈时间}：作者提交后多久收到第一批有效评论
    \item \textbf{缺陷逃逸率}：被合并后仍流入生产环境的问题比例
    \item \textbf{假阳性率}：AI 评论中被证明无效的比例
    \item \textbf{重开 PR 比例}：说明前期审查是否真正发现核心问题
    \item \textbf{reviewer 负担}：资深工程师在低价值评论上消耗了多少时间
\end{itemize}

\begin{knowledgebox}{好的审查系统不是让每个 PR 更吵，而是让高风险 PR 更清晰}
如果 AI 给每个 diff 都生成大量模板化评论，团队很快就会形成忽视习惯。真正有效的系统应该把注意力集中在最关键的问题上，让重要评论更容易被认真处理。
\end{knowledgebox}

\subsection{本章小结}

AI 代码审查工具通过自动化常规检查显著提升了效率和一致性，但复杂的业务逻辑、架构决策和仓库特有的最佳实践仍然需要人类审查者。在 AI 编码时代，代码审查的重要性不降反升。

